\documentclass[11pt]{article}

\usepackage[T1]{fontenc}
\usepackage[utf8]{inputenc}
\usepackage[margin=1in]{geometry}
\usepackage{amsmath,amssymb,amsthm}
\usepackage[shortlabels]{enumitem}
\usepackage{microtype}

\setlength{\parindent}{0pt}
\setlength{\parskip}{0.65em}
\setlist[enumerate]{leftmargin=*,itemsep=0.45em,topsep=0.35em}

\newcommand{\N}{\mathbb{N}}
\newcommand{\Z}{\mathbb{Z}}
\newcommand{\Q}{\mathbb{Q}}
\newcommand{\R}{\mathbb{R}}
\newcommand{\problem}[2]{\subsection*{Problem #1 \hfill \normalfont[#2 points]}}

\begin{document}
\pagestyle{empty}

\begin{center}
    {\Large\bfseries 21-128 and 15-151}\\[3pt]
    {\LARGE\bfseries Problem Sheet 2}\\[7pt]

Solutions to the following six exercises and optional bonus problem are to be 
submitted to Gradescope by

\textbf{11 PM on Wednesday, the 9th of September, 2026.}

\end{center}



\hrule
\vspace{0.6em}

\textbf{Instructions.}
This assignment has 30 points, together with an optional 1-point bonus problem.
Unless a problem explicitly says otherwise, every answer must be justified by a proof.
A counterexample must satisfy all of the hypotheses and must be checked explicitly.
We use $\subseteq$ for containment and $\subsetneq$ for proper containment.

\problem{1}{4}
For each statement below, determine whether it is true for every choice of nonempty sets $A,B,C\subseteq\Z$.
Prove each statement that is always true, and give a fully specified counterexample to each statement that is not always true.
\begin{enumerate}[(a)]
    \item
    \[
        A\cap B\cap C\subsetneq A\cup B.
    \]

    \item
    \[
        (A\setminus B)\cap(B\setminus A)=\varnothing.
    \]

    \item
    \[
        \bigl(A\cap B\ne\varnothing\bigr)
        \Longrightarrow
        \bigl(A\setminus B\subsetneq A\bigr).
    \]
\end{enumerate}

\problem{2}{5}
For each $n\in\Z_{>0}$, define
\[
    M_n=\{x\in\Z:n\mid x\}.
\]
Let $m,n\in\Z_{>0}$.
\begin{enumerate}[(a)]
    \item Prove that
    \[
        M_m\subseteq M_n
        \quad\Longleftrightarrow\quad
        n\mid m.
    \]

    \item Prove that
    \[
        M_m=M_n
        \quad\Longleftrightarrow\quad
        m=n.
    \]
\end{enumerate}

\problem{3}{4}
For sets $X$ and $Y$, their \emph{symmetric difference} is
\[
    X\mathbin{\triangle}Y=(X\setminus Y)\cup(Y\setminus X).
\]
Prove that, for all sets $A,B,C$,
\[
    C\setminus(A\mathbin{\triangle}B)
    =
    (A\cap B\cap C)\cup\bigl(C\setminus(A\cup B)\bigr).
\]
Your proof must proceed by showing that an arbitrary element belongs to the left-hand side if and only if it belongs to the right-hand side. A Venn diagram may be used for discovery, but it is not a proof.

\problem{4}{5}
A set $T$ is called \emph{transitive} if, for all sets $x$ and $y$,
\[
    (x\in y\ \wedge\ y\in T)\Longrightarrow x\in T.
\]
\begin{enumerate}[(a)]
    \item Determine, with proof, which of the following sets are transitive:
    \[
        T_1=\{\varnothing,\{\varnothing\}\},
        \qquad
        T_2=\{\{\varnothing\}\}.
    \]

    \item Prove that if $T$ is transitive, then its power set $\mathcal{P}(T)$ is also transitive.
    Your proof should explicitly distinguish membership from containment.
\end{enumerate}

\problem{5}{5}
A set $S\subseteq\R$ is called \emph{open} if
\[
    (\forall x\in S)(\exists\varepsilon>0)(\forall y\in\R)
    \bigl(|x-y|<\varepsilon\Longrightarrow y\in S\bigr).
\]
Let $S\subseteq\R$ be open, and suppose that $a,b\in S$ with $a<b$.
Prove that $S\cap(a,b)$ contains infinitely many distinct points.
Your proof should explicitly construct such points from an $\varepsilon$ supplied by the definition of openness.

\problem{6}{7}
Define
\[
    f:\Z_{>0}\times\Z_{>0}\longrightarrow\R,
    \qquad
    f(a,b)=\frac{(a+1)(a+2b)}{2}.
\]
\begin{enumerate}[(a)]
    \item Prove that the image of $f$ is contained in $\Z_{>0}$.

    \item Determine
    \[
        \operatorname{im}(f)\cap\{1,2,\ldots,20\},
    \]
    and use your answer to state a conjecture describing the entire image of $f$.

    \item Prove your conjecture: determine exactly which positive integers belong to the image of $f$.
\end{enumerate}

\textit{Hint for part (c).}
Set $u=a+1$ and $v=a+2b$. Observe that $2f(a,b)=uv$, that $2\le u<v$, and that $u$ and $v$ have opposite parity. For the reverse direction, it may be useful to write a positive integer as $2^kq$ with $q$ odd.

\newpage
\subsection*{Bonus Problem \hfill \normalfont[1 bonus point]}
A function $u:\R\to\R$ is \emph{even} if $u(-x)=u(x)$ for every $x\in\R$, and is \emph{odd} if $u(-x)=-u(x)$ for every $x\in\R$.
Let $f:\R\to\R$ be an arbitrary function.
\begin{enumerate}[(a)]
    \item Prove that there exists a unique pair of functions $g,h:\R\to\R$ such that $g$ is even, $h$ is odd, and
    \[
        f=g+h.
    \]

    \item Suppose that
    \[
        f(x)=c_0+c_1x+\cdots+c_nx^n
    \]
    is a polynomial function. Express the functions $g$ and $h$ from part (a) directly in terms of the coefficients $c_0,c_1,\ldots,c_n$, and prove that your formulas work.
\end{enumerate}

\end{document}
